\color{unidarkgreen}{\subsection{Автоматическое повторное включение}}
\color{black}

\begin{enumerate}[label=\arabic{section}.\arabic{subsection}.\arabic*,  align=left, labelsep=4pt, leftmargin=0pt, itemindent=57pt]

\item
Автоматическое повторное включение (АПВ) предназначено для быстрого восстановления
первоначального состояния электрической сети после аварийного отключения путём автоматического
включения выключателя.

\vspace{2mm}
\begin{figure}[H]
\centering
\includegraphics[width=0.9\textwidth,height=0.9\textheight,keepaspectratio]{\fbpath/661.3011 АПВ/img/APV_ShSV.pdf}
\captionof{figure}{Функциональная схема логики АПВ}\label{pict_apv:APV_ShSV}
\end{figure}

\item
Выбор действующих режимов АПВ реализуется тремя командами управления режимами АПВ:
\begin{itemize}
    \item \textit{«АПВ ОНШ1 + ННШ2»} (включение с контролем отсутствия напряжения на 1СШ и наличия напряжения на 2СШ) -- АПВ 1СШ;
    \item \textit{«АПВ ОНШ2» + ННШ1»} (включение с контролем отсутствия напряжения на 2СШ и наличия напряжения на 1СШ) -- АПВ 2СШ;
    \item \textit{«АПВ без контролей»} (включение без контроля напряжений).
\end{itemize}

\item
Возможные режимы АПВ приведены в таблице \ref{apv:tbl1}.

\small
\renewcommand{\arraystretch}{1}
\setlength{\extrarowheight}{2pt}

\begin{longtable}{|>{\centering\arraybackslash}m{5.5cm}|>{\centering\arraybackslash}m{3.5cm}|>{\centering\arraybackslash}m{3.5cm}|>{\centering\arraybackslash}m{3.5cm}|}
\caption{Условия включения\hfill\vspace{-0.5\baselineskip}}\label{apv:tbl1}\\
\hline
\rowcolor{gray!30}
& \multicolumn{3}{c|}{Команды управления режимами АПВ} \\
\hhline{>{\arrayrulecolor{gray!30}}->{\arrayrulecolor{black}}|-|-|-|}
\rowcolor{gray!30}
\raisebox{3pt}{\makecell[c]{Режим АПВ}} & 
\makecell[c]{АПВ ОНШ1+ННШ2} & 
\makecell[c]{АПВ ОНШ2+ННШ1} & 
\makecell[c]{АПВ без контролей} \\
\hline
\endfirsthead

\caption*{\hspace{3pt}\emph{Продолжение таблицы \ref{apv:tbl1}}\hfill\ }\vspace{-0.5\baselineskip} \\
\hline
\rowcolor{gray!30}
& \multicolumn{3}{c|}{Входные сигналы ФСУ} \\
\hhline{>{\arrayrulecolor{gray!30}}->{\arrayrulecolor{black}}|-|-|-|}
\rowcolor{gray!30}
\raisebox{3pt}{\makecell[c]{Режим АПВ}} & 
\makecell[c]{АПВ ОНШ1+ННШ2} & 
\makecell[c]{АПВ ОНШ2+ННШ1} & 
\makecell[c]{АПВ без контролей} \\
\hline
\endhead

\hline
\makecell[c]{КС(УС)} & 0 & 0 & 0 \\
\hline
\makecell[c]{ОНШ1+ННШ2 / КС(УС)} & 1 & 0 & 0 \\
\hline
\makecell[c]{ОНШ2+ННШ1 / КС(УС)} & 0 & 1 & 0 \\
\hline
\makecell[c]{ОНШ1+ННШ2 / ОНШ2+ННШ1 /  КС(УС)} & 1 & 1 & 0 \\
\hline
\makecell[c]{Без контролей} & 0/1 & 0/1 & 1 \\
\hline
\end{longtable}
\normalsize

\item
Как видно из таблицы \ref{apv:tbl1}, режим АПВ с КС(УС) введен всегда и выводится посредством входного сигнала ФСУ \textit{«АПВ
без контролей»}. А режимы ОНШ1+ННШ2 и ОНШ2+ННШ1  могут быть введены одновременно. Условия включения
в различных режимах более подробно описаны {\color{red}в пункте 14} настоящего руководства по эксплуатации.

\item
Пуск АПВ линии осуществляется по цепи несоответствия при аварийном отключении выключателя
(по наличию входного сигнала ФСУ \textit{«Сигнал несоотв. В»}). Срабатывание происходит при набранной выдержке
времени готовности АПВ с контролем заданного режима и выполнения условий включения. Выдержка времени
готовности АПВ после оперативного включения и первого цикла регулируется уставкой \textbf{«Тгот АПВ»}.

\item
Ввод в работу АПВ осуществляется с помощью программного ключа \textbf{«АПВ»}. Цикл АПВ происходит при
выполнении условий включения с выдержкой времени \textbf{«Тср АПВ»}. В случае сброса сигнала пуска АПВ при
отсчете времени \textbf{«Тср АПВ»} цепь логики сбрасывается, и схема АПВ возвращается в исходное состояние
(готовность АПВ с набранной выдержкой готовности). Длительность импульса включения выключателя от АПВ
задается уставкой \textbf{«Тимп АПВ»}.

\item
Предусмотрено ограничение ожидания условий включения от АПВ, которое вводится программным
ключом \textbf{«Огр ожид усл вкл В»}. Если в течение выдержки времени \textbf{«Тожид усл вкл В»} после аварийного
отключения не возникают условия включения, то действие АПВ запрещается.

\item
Оперативный ввод/вывод АПВ производится с помощью БУ (блоков управления функциями) \textit{«Вывод АПВ»}. 

\item
Кроме того, выставлением соответствующих уставок может осуществляться запрет АПВ при
срабатывании ступеней ДЗ-ФЗ и ДЗ-ФФ, ступеней ТНЗНП, ступеней МТЗ, ступеней АМТЗ, при срабатывании
оперативного ускорения (как с выдержкой, так и без выдержки времени) ДЗ и ТНЗНП, автоматического
ускорения, АОДС и внешнего отключения. На запрет АПВ также жёстко действуют срабатывание УРОВ, ЗНФ,
ЗНР, оперативное отключение, режим местного управление от привода В, отключение от аварийного снижения
давления элегаза в ТТ. Также для запрета АПВ предусмотрены входные сигналы ФСУ \textit{«Внешний запрет АПВ»}.

\vspace{2mm}
\begin{figure}[H]
\centering
\includegraphics[width=0.9\textwidth,height=0.9\textheight,keepaspectratio]{\fbpath/661.3011 АПВ/img/APV_Logic.pdf}
\captionof{figure}{Функциональная схема запрета АПВ}\label{pict_apv:APV_Logic}
\end{figure}


\small
\begin{longtable}{|>{\centering\arraybackslash}m{5.3cm}|>{\centering\arraybackslash}m{3.3cm}|>{\centering\arraybackslash}m{4.2cm}|>{\centering\arraybackslash}m{1.8cm}|>{\centering\arraybackslash}m{1cm}|}
\caption{Перечень уставочных параметров АПВ \hfill\vspace{-0.5\baselineskip}}\label{desc_apv:tbl2}\\ 
\hline
\rowcolor{gray!30}
Параметр & Условное обозначение на схеме & Значение/ Диапазон & Единица измерения & Шаг \\ 
\hline
\endfirsthead
\caption*{\hspace{3pt}\emph{Продолжение таблицы \ref{desc_apv:tbl2}\hfill\vspace{-0.5\baselineskip}}} \\ % сделано по ГОСТ 2.105 п.6.8.7
\hline
\rowcolor{gray!30}
Параметр & Условное обозначение на схеме & Значение/ Диапазон & Единица измерения & Шаг \\ 
%\hline
\endhead

%===m>АПВ |-
%===m
\end{longtable}
\normalsize





















\end{enumerate}